\documentclass[12pt]{article}
\usepackage[utf8]{inputenc}
\usepackage[T1]{fontenc}
\usepackage[margin=1in]{geometry}
\usepackage{setspace}
\usepackage{xcolor}
\usepackage[colorlinks=true, urlcolor=blue!60!black, linkcolor=blue!60!black]{hyperref}
\setlength{\parindent}{0em}
\setlength{\parskip}{0.45em}
\renewcommand{\baselinestretch}{1.0}
\begin{document}
\begin{center}
{\Large \textit{Hicklin v. Orbeck}, 437 U.S.\ 518 (1978)\par}
\vspace{0.3em}
{\normalsize Case Brief --- One-Page Edition\par}
\end{center}

\textbf{Memory Jogger.} ``Alaska Hire'' fell under Article IV: resource ownership did not let Alaska make private employers prefer residents (\href{https://supreme.justia.com/cases/federal/us/437/518/}{437 U.S.\ 518}).

\textbf{Facts.} Alaska required state oil and gas leases and pipeline permits to mandate hiring residents first, reaching contractors, suppliers, and refineries. When enforcement began at the peak of Trans-Alaska Pipeline construction, Sidney Hicklin and other workers who could not obtain resident cards were barred from pipeline jobs.

\textbf{Procedural History.} The Alaska Supreme Court struck the one-year residency rule but, 3--2, upheld the preference based on state ownership of oil and gas.

\textbf{Votes.} Reversed, 9--0, by Justice Brennan for a unanimous Court.

\textbf{Holding.} Discrimination against nonresidents in private employment requires that they be a ``peculiar source of the evil'' and a close fit; state ownership is a factor, not an exemption.

\textbf{Reasoning.} Alaska's unemployment stemmed from lack of training and remoteness, not nonresidents, and the flat preference aided employed and unemployed residents alike. The Act reached employers with no tie to state land or funds, and the oil was bound for interstate commerce.

\textbf{Examples.} \emph{Same-Side:} a wind-farm supplier hiring preference. \emph{Opposite-Side:} higher nonresident hunting fees. \emph{Fence-Sitter:} a city's resident quota on city-funded projects.

\textbf{Critique.} Blending Article IV with Commerce Clause cases left the needed proprietary nexus unclear, and States have few tools to steer resource-boom jobs to undertrained local residents.

\textbf{Key Quotation.} ``In sum, the Act is an attempt to force virtually all businesses that benefit in some way from the economic ripple effect of Alaska's decision to develop its oil and gas resources to bias their employment practices in favor of the State's residents.'' --- Justice Brennan, 437 U.S.\ at 531.

\end{document}
